% Chapter 6: Boolean Algebra over Trapdoor Values
% DEFINES: ch:boolean-algebra; sec:bit-images, sec:union-intersection,
%   sec:complement-approximate, sec:equality-channel, sec:boolean-chains,
%   sec:ch6-notes; def:bit-image, thm:not-approximate, ex:boolean-watchlist,
%   rem:equality-channel, tab:boolean-ops
% REFERENCES (backward): sec:strategies, def:correctness, thm:composition,
%   def:uniformity, ex:watchlist-numbers
% REFERENCES (forward, part labels): part:confidentiality, part:constructions

\chapter{Boolean Algebra over Trapdoor Values}
\label{ch:boolean-algebra}

Part~II built the cipher map and its parameters; this part asks what
algebra those objects support. The most concrete answer comes first. A set
can be turned into a single trapdoor value, a bit string, on which the
Boolean operations, union, intersection, and complement, are computed
directly by bit operations, with no decoding. The result is an online
cipher construction in the sense of \Cref{sec:strategies}: nothing is
pre-solved, and the operations are exactly those the bit representation
supports. It is also where a structural fact of the whole framework first
appears: union and intersection pass through cleanly, but complement cannot,
and the reason is a counting argument that no construction escapes.

\section{Sets as Bit Images}
\label{sec:bit-images}

Fix a hash $h$ that sends each element of the universe to a position in an
$n$-bit string. A set becomes the union of its elements' bits.

\begin{definition}[Bit image]
\label{def:bit-image}
For a set $A$ drawn from the universe, the \emph{bit image} is
\[
  F(A) \;=\; \bigvee_{a \in A} h(a) \;\in\; \B^n,
\]
the bitwise OR of the one-bit hashes of its elements. Membership is tested
by the bits: $x$ is reported in $A$ when $h(x)$'s bit is set in $F(A)$.
\end{definition}

The image $F(A)$ is a cipher value for the set $A$: the untrusted machine
holds the bit string and can compute with it, but without $h$ it cannot
read off which elements are present. Membership testing is one-sided. If
$x \in A$ its bit is set, so a true member always tests positive; but the
bit may also have been set by some other element $y \in A$ with
$h(y) = h(x)$, a collision, in which case a non-member tests positive too.
Those collisions are exactly the false positives, the correctness error
$\eta$ of \Cref{def:correctness}, here realized as a Bloom-filter-style
membership structure. The construction is deterministic: a set has one
image, so $K(x) = 1$ in the language of \Cref{def:uniformity}. That
determinism is what makes the algebra below exact, and, as
\Cref{sec:equality-channel} shows, what it costs.

\section{Union and Intersection}
\label{sec:union-intersection}

The two binary operations behave differently, and the difference is the
chapter's first theme.

Union is exact. Because $F$ is an OR of element bits, and OR distributes
over the set union,
\[
  F(A \cup B) \;=\; F(A) \mathbin{|} F(B),
\]
with no error: every bit set by an element of $A \cup B$ is set on one side
or the other, and none is spuriously introduced. Intersection is only
one-sided. The bitwise AND keeps the positions set in both images, but a
position can be set in both for the wrong reason, an element of
$A \setminus B$ and an element of $B \setminus A$ that happen to hash to the
same bit. Such cross-collisions survive the AND as spurious members, so
\[
  F(A \cap B) \;\subseteq\; F(A) \mathbin{\&} F(B),
\]
with equality exactly when no cross-collision occurs. The spurious bits are
false positives for the intersection, never missing bits.
\Cref{tab:boolean-ops} collects the operations.

\begin{table}[t]
\centering
\caption{The set operations as bit operations, and whether each is exact.}
\label{tab:boolean-ops}
\begin{tabular}{@{}lll@{}}
\toprule
Set operation & Bit operation & Exact? \\
\midrule
$A \cup B$ & $F(A) \mathbin{|} F(B)$ & exact \\
$A \cap B$ & $F(A) \mathbin{\&} F(B)$ & one-sided (spurious bits from collisions) \\
$A^{\complement}$ & ${\sim} F(A)$ & approximate (\Cref{thm:not-approximate}) \\
$\varnothing$ & $0^n$ & exact \\
\bottomrule
\end{tabular}
\end{table}

\begin{example}[Combining two watchlists]
\label{ex:boolean-watchlist}
Take $n = 8$ (here the bit-image width, distinct from the $12$-bit token
space of \Cref{ex:watchlist-numbers}) and the watchlist
$S = \{\mathrm{alice}, \mathrm{carol}\}$ beside a second list
$V = \{\mathrm{carol}, \mathrm{dave}\}$. Numbering bit positions from the
right ($0$ least significant), suppose $h$ places alice at bit $1$, carol at
bit $6$, and dave at bit $2$, so
\[
  F(S) = \mathtt{0100\,0010}, \qquad F(V) = \mathtt{0100\,0100} .
\]
Their union is exact: $F(S) \mathbin{|} F(V) = \mathtt{0100\,0110}$, the bits
for alice, carol, and dave, which is precisely $F(S \cup V)$ for
$S \cup V = \{\mathrm{alice}, \mathrm{carol}, \mathrm{dave}\}$. Their
intersection here is also clean, $F(S) \mathbin{\&} F(V) =
\mathtt{0100\,0000} = F(\{\mathrm{carol}\}) = F(S \cap V)$, because alice and
dave land on different bits. Had they collided, say both on bit $2$, the AND
would have kept bit $2$ as well, reporting a phantom member of $S \cap V$;
that is the one-sided error the inclusion above allows.
\end{example}

\section{The Complement Is Approximate}
\label{sec:complement-approximate}

Union and intersection cost nothing or cost only spurious bits. Complement
is different in kind: there is no exact bit operation for it, and the
obstruction is a counting argument, not a defect of this particular hash.

\begin{theorem}[The complement asymmetry]
\label{thm:not-approximate}
Let the universe be larger than the bit width, so that every bit position is
set by some element. Then the bitwise complement ${\sim}F(A)$ satisfies
\[
  {\sim}F(A) \;\subseteq\; F(A^{\complement}),
\]
and the inclusion is strict whenever some bit is set by both an element of
$A$ and an element of $A^{\complement}$. The gap grows with $\lvert A\rvert$:
the approximation is tightest when $A$ is small and degrades as $A$ fills
the bit space.
\end{theorem}

\begin{proof}
A bit lies in ${\sim}F(A)$ exactly when no element of $A$ sets it. Since the
universe is larger than the bit width, every bit is set by some element of
the universe $A \cup A^{\complement}$; a bit set by no element of $A$ must
then be set by some element of $A^{\complement}$, hence lies in
$F(A^{\complement})$. So ${\sim}F(A) \subseteq F(A^{\complement})$. The
reverse fails for any bit set both by an element of $A$ and by an element of
$A^{\complement}$: such a bit is in $F(A^{\complement})$ but, being set in
$F(A)$, not in ${\sim}F(A)$. The number of such doubly-set bits increases
with $\lvert A\rvert$, since a larger $A$ sets more of the bit space and so
shares more positions with $A^{\complement}$. When $A$ is small,
$A^{\complement}$ is large and saturates the bits, $F(A^{\complement})
\to 1^n$, while ${\sim}F(A)$, complementing a sparse $F(A)$, is also near
$1^n$, and the two agree.
\end{proof}

The direction matters, and an early account of this construction got it
backwards: complement is a \emph{better} approximation for \emph{small} sets
and a worse one as the set grows, not the reverse. Intuitively, a small set
leaves most of the bit space unset, and its complement genuinely fills that
space; a large set has already crowded the bits, and complementing its image
loses exactly the positions it shares with everything outside it. The error
is governed by the ratio $\lvert A\rvert / 2^n$, the same fill factor that
controls collisions, and it is structural: no finite-width construction can
make complement exact over an unbounded universe, because the universe
cannot be injected into $\B^n$. This asymmetry, exact joins, approximate
complement, is inherited by every system built on a bitwise trapdoor
algebra, and it returns whenever such a system is analyzed.

\section{The Equality Channel}
\label{sec:equality-channel}

The determinism that makes union and intersection exact has a price, and
naming it explains why this online algebra trades away confidentiality for
expressiveness. Because $K = 1$, a set has a single image, so two sets are
equal exactly when their images are: $F(A) = F(B)$ if and only if $A = B$
(up to the collision error). An untrusted machine that holds two cipher
sets can therefore test them for equality, and more generally can read the
equality structure of everything it is given, which queries repeat, which
results coincide, even though it cannot read the contents. This is the
\emph{equality channel}, the leakage characteristic of a deterministic
trapdoor representation.

\begin{remark}[The cost of determinism]
\label{rem:equality-channel}
The equality channel is the flip side of \Cref{def:uniformity}. Hiding
\emph{which} value a token encodes needed multiple representations,
$K(x) > 1$, so that one value spreads across many tokens and equal values no
longer produce equal tokens. The bit algebra here is the opposite extreme,
$K = 1$: a single deterministic image per set, which is what lets the
operations combine images directly, but which leaks equality in full. The
two cannot both be had in this construction. Randomizing the image to raise
$K$ would break the bit operations, since two random images of the same set
no longer OR or AND meaningfully. The online Boolean algebra is therefore
the exact-operations, equality-leaking corner of the design space; buying
confidentiality back means leaving it for a batch construction with
$K > 1$, at the cost of the clean algebra.
\end{remark}

\section{Boolean Chains and Their Error}
\label{sec:boolean-chains}

A Boolean circuit over cipher sets is a composition of cipher maps, so its
error is governed by \Cref{thm:composition}, with one addition. Through a
chain of joins and meets, the only error is the spurious-bit error of
intersection, and it accrues exactly as the composition theorem predicts:
the chain's correctness is bounded by $1 - \prod_i (1 - \eta_i)$, the
inequality of \Cref{sec:composability}, with each $\eta_i$ the collision
error of one operation. A circuit of unions and bounded intersections stays
predictable in the way \Cref{ch:four-properties} promised.

Complement breaks the pattern. Each NOT contributes not the tame collision
error but the structural error of \Cref{thm:not-approximate}, which scales
with the fill factor $\lvert A\rvert / 2^n$ of its operand rather than with a
per-operation budget. A circuit dense in complements therefore degrades
faster than one built from joins and meets alone, and the degradation is not
removed by enlarging the structure, only deferred. The practical reading is
the one the asymmetry forces: design Boolean cipher programs to lead with
unions and intersections and to use complement sparingly and on small sets,
where \Cref{thm:not-approximate} is tight.

\section{Notes and Provenance}
\label{sec:ch6-notes}

The bit-image algebra is the trapdoor Boolean construction of the
foundations notes, and the deterministic $K = 1$ baseline is the kernel of
the retired boolean-algebra-over-trapdoor-sets paper, folded here. The
operation table and the exact/approximate split follow the spine's
treatment. The complement asymmetry (\Cref{thm:not-approximate}) is stated
in its corrected direction: the original blog account claimed complement
improves for large sets, which is backwards, and the spine carries the
erratum; this chapter proves the correct direction from the pigeonhole
argument. The equality channel is the leakage regime of the deterministic
construction, the price the online algebra pays for exact operations, and it
is the concrete motivation for the multiple-representation machinery of
\Cref{def:uniformity}. The full Boolean cipher type, its noise behavior, and
the experimental validation of the asymmetry are developed in
\Cref{part:constructions}; what the equality channel leaks is the cost of
determinism, and the multiple representations of \Cref{def:uniformity} are
what buy it back, while at the multi-instance scale it reappears as the
coincidence leakage that \Cref{part:confidentiality} measures.
